\section{Step 76 Load neighbors from reconstructed samples}
Fill \texttt{row\_neighbors}. Load the corresponding upper sample into \texttt{\$t4}. Initialize left and upper-left to zero in \texttt{\$t5} and \texttt{\$t6}. Determine the current sample’s offset from \texttt{\$a0}, the row’s first sample. For offsets below three, keep the two left values zero. Otherwise load left from destination minus three and upper-left from upper cursor minus three.

The distance is three bytes because a red sample depends on red in the preceding pixel, green on green, and blue on blue. A load at minus one would connect different channels. The left load observes a sample already reconstructed by this loop; using its original residual would create a different recurrence.

Branch to the supplied \texttt{row\_neighbors\_ready} checkpoint for the first three samples. For the first byte of the worked Up row, the neighbors are upper 10, left zero and upper-left zero. At later columns, all three registers contain unsigned byte values. Retain the full-width cursors; only the loaded sample values are eight bits.

\paragraph{Observe} Use \texttt{s76-1}; stop before \texttt{row\_neighbors\_ready}. Inspect the state described above before continuing.

\Needspace{7\baselineskip}
\paragraph{Check} What is the left-neighbor distance in RGB8?
\begin{itemize}
\item[A.] One byte.
\item[B.] Three bytes.
\item[C.] One complete scanline.
\end{itemize}
